\section{Conclusiones}
En este trabajo de revisión hemos resuelto analíticamente la dinámica vibracional y las reglas de selección óptica de dos sistemas moleculares arquetípicos: el benceno plano ($D_{6h}$, $N=12$) y el buckminsterfullereno esférico ($\text{C}_{60}$, $Y_h$, $N=60$). A partir de los teoremas de átomos no desplazados formulados por Ramadevi y Dubey~\cite{Ramadevi2019}, se redujeron las representaciones mecánicas $\Gamma_{3N}$ aislando traslaciones y rotaciones rígidas, deduciendo de forma exacta los 30 modos normales del benceno y los 174 modos del $\text{C}_{60}$, los cuales colapsan en 46 frecuencias propias fundamentales debido a la alta degeneración espacial impuesta por la simetría icosaédrica.

El análisis de paridad fundamentó rigurosamente la regla de exclusión mutua en ambos grupos centrosimétricos ($G \cong H \times C_i$): los operadores dipolares de absorción infrarroja (sector impar $u$) y el tensor de polarizabilidad de dispersión Raman (sector par $g$) son mutuamente excluyentes, confinando al régimen ópticamente silente a 12 modos en el benceno ($40\%$) y a 120 modos en el $\text{C}_{60}$ ($69\%$). Asimismo, el producto tensorial simétrico $(T_{1u} \otimes T_{1u})_s = A_g \oplus H_g$ demostró que los 8 modos cuadrupolares $H_g$ rigen tanto la inestabilidad dinámica de Jahn--Teller en el radical $\text{C}_{60}^-$ como el emparejamiento de Cooper en sales superconductoras $A_3\text{C}_{60}$. Finalmente, para superar la inviabilidad tipográfica de imprimir las 120 matrices de orden $3 \times 3$ de $Y_h$, se implementó un simulador 3D interactivo en WebGL de acceso abierto, permitiendo la exploración geométrica y dinámica completa de estos grupos.
